% !TEX program = pdflatex
% ---------------------------------------------------------------------------
% Implementation notes for the numerics of
%   "Finite de Finetti for convex bodies and polynomial optimization"
% Standalone companion document for the GitHub repository.
% Compile with:  pdflatex implementation_notes.tex  (twice)
% ---------------------------------------------------------------------------
\documentclass[11pt,a4paper]{article}

\usepackage[T1]{fontenc}
\usepackage[utf8]{inputenc}
\usepackage{geometry}
\geometry{a4paper, total={150mm,225mm}}
\usepackage{amsmath,amssymb,amsthm,mathtools}
\usepackage{array}
\usepackage{booktabs}
\usepackage{enumitem}
\usepackage[dvipsnames]{xcolor}
\usepackage[colorlinks, citecolor=PineGreen, linkcolor=MidnightBlue, urlcolor=PineGreen]{hyperref}
\usepackage{xurl}

% ------------------------------------------------------------- paper macros
\newcommand*{\tmin}{\dot{\otimes}}
\newcommand*{\tmax}{\hat{\otimes}}
\newcommand*{\lrbracket}[1]{\left(#1\right)}
\newcommand*{\lrvert}[1]{\left\lvert #1 \right\rvert}
\newcommand*{\Acal}{\mathcal{A}}
\newcommand*{\Bcal}{\mathcal{B}}
\newcommand*{\Xcal}{\mathcal{X}}
\newcommand*{\Ycal}{\mathcal{Y}}
\DeclareMathOperator{\ext}{ext}
\DeclareMathOperator{\Val}{Val}
\DeclareMathOperator{\conv}{conv}
\newcommand*{\Kasm}{K_{\mathrm{asm}}}
\newcommand*{\Kmm}{K_{\mathrm{mm}}}
\newcommand*{\Casm}{C_{\mathrm{asm}}}
\newcommand*{\Cmm}{C_{\mathrm{mm}}}
% one place to fill in the final repository location:
\newcommand*{\repourl}{\url{https://github.com/JuliusAZeiss/Finite-de-Finetti-for-convex-bodies-and-Polynomial-Optimization}}
\newcommand*{\code}[1]{\texttt{#1}}
% allow a little slack around long typewriter file names
\emergencystretch=2em

\theoremstyle{remark}
\newtheorem*{note}{Note}

\title{Implementation notes and reproduction guide\\[2mm]
{\large Numerics companion to
``Finite de Finetti for convex bodies and polynomial optimization''~\cite{paper}}}
\author{Julius A. Zeiss \and Gereon Kossmann \and Ren\'e Schwonnek \and Martin Pl\'avala}
\date{August 2026}

\begin{document}
\maketitle

\noindent
This document collects the implementation details behind the section
\emph{Numerical illustration: non-local games over polygon theories} of the
paper~\cite{paper}, which the paper itself states only in compressed form.
It is meant to be read alongside the code in the
repository\footnote{\repourl\\ File names in typewriter font refer to
files in the repository; \code{unreduced\_checks/} is its only subfolder.}
and should let a reader check, modify and re-run every reported number
without reverse-engineering the source. No result stated here is new;
everything expands on the numerics section of the paper, in particular on
its Table~1, its Figure~1 and its observations (a)--(f).

\tableofcontents

% ===========================================================================
\section{What is computed}\label{sec:overview}

For four polytopic experiments --- the classical bit, the square gbit, the
strongly self-dual regular pentagon, and the pentagon with an additional
local inequality constraint --- the code computes three things for the
assemblage value $\gamma^*(G)$ of the CHSH game:
\begin{enumerate}[leftmargin=2em]
\item the outer levels $p^{(n)}$, $n=1,\ldots,5$, of the convergent
hierarchy of the paper (theorem \emph{Convergence of the hierarchy}; label
\code{thm:convergence\_hierarchy} in the source), each solved as an
\emph{exact} linear program (Sections~\ref{sec:lp}--\ref{sec:reductions});
\item at every level a feasible product strategy with value $L^{(n)}$,
obtained by the rounding scheme of the theorem \emph{Convergent inner
sequence} (label \code{thm:inner\_bounds}) from the optimizer returned by
the LP solver (Section~\ref{sec:rounding});
\item the exact optimum $\gamma^*(G)$ itself, computed directly --- bypassing
hierarchy and rounding --- and certified in exact arithmetic
(Section~\ref{sec:exact}), as the reference value of the table.
\end{enumerate}
Together these sandwich the optimum, $L^{(n)}\leq\gamma^*(G)\leq p^{(n)}$,
with explicit certificates on both sides at every level. The remaining
sections document the problem sizes and runtimes (Section~\ref{sec:sizes}),
an independent re-computation of all outer levels \emph{without} any of the
reductions of Section~\ref{sec:reductions} (Section~\ref{sec:unreduced}),
the study behind observation (e) of the paper --- the dependence of the
rounded values on the chosen optimizer (Section~\ref{sec:sensitivity}) ---,
the level-1 diagnostics behind observation (a) (Section~\ref{sec:level1}),
and finally how to reproduce everything and what each file contains
(Sections~\ref{sec:repro}--\ref{sec:files}).

Table~\ref{tab:map} maps the claims of the paper's numerics section to the
files that produce them.

\begin{table}[ht]\centering\small
\begin{tabular}{>{\raggedright\arraybackslash}p{56mm}>{\raggedright\arraybackslash}p{84mm}}
\toprule
claim in the paper & produced / checked by\\
\midrule
Table 1, rows $p^{(n)}$ and $L^{(n)}$; Figure 1 & \code{reproduce.py} $\to$ \code{results.json}; \code{make\_figure.py}\\
Table 1, column $\gamma^*$ (exact) & \code{exact\_game\_value.py} $\to$ \code{exact\_game\_values.json}\\
(a) level 1 is the maximal tensor product & \code{level1\_pseudocorrelations.py}\\
(b) zero-gap certificates (bit $n=2$, gbit $n=4$) & \code{results.json}\\
(c) strict convergence, $\approx 0.1/n$ & \code{results.json}, \code{make\_figure.py}\\
(d) certified inner strategies, bracket at $n=5$, $m^*=n-1$ & \code{results.json} (field \code{min\_slack}, \code{per\_m\_best})\\
(e) optimizer dependence of $L^{(n)}$, quoted ranges & \code{unreduced\_checks/direct\_inner\_sensitivity.py}\\
(f) the inequality constraint binds & \code{exact\_game\_values.json}, \code{results.json}\\
``every level is an exact LP'' (the reductions are exact) & \code{unreduced\_checks/direct\_hierarchy.py}, \code{unreduced\_checks/direct\_hierarchy\_tree.py}\\
all of the above in one executable document & \code{reproduce\_results.ipynb}\\
\bottomrule
\end{tabular}
\caption{Where each claim of the numerics section is produced.}\label{tab:map}
\end{table}

% ===========================================================================
\section{The optimization problem}\label{sec:problem}

Throughout, $G$ is the CHSH game:
$\Acal=\Bcal=\Xcal=\Ycal=\lbrace 0,1\rbrace$, uniform question
distribution $\pi$, and predicate $V(a,b,x,y)=1$ iff $a\oplus b=x\wedge y$;
the game coefficients are $g_{abxy}=\pi(x,y)V(a,b,x,y)$ and all values are
reported in the winning-probability convention. The assemblage value of the
game is the bilinear problem
\begin{align}\label{eq:gamma}
\begin{split}
    \gamma^*(G) \;=\; \sup_{\xi_B,\, N_B}\quad
    & G_{\Acal\Bcal\Xcal\Ycal}\!\left[\lrbracket{N_B\otimes \xi_B}\lrbracket{\chi_{A(K_B)^{\star}}}\right]\\
    \text{s.t.}\quad
    & \xi_B \in \Kasm \coloneqq CS^1_{2,2}\tmax K_B,\\
    & N_B \in \Kmm \coloneqq \left\lbrace N\in A(K_B)^+\tmax CS^+_{2,2} \,:\,
      \lrbracket{\operatorname{id}\otimes 1_{CS^1_{2,2}}}\lrbracket{N}=1_{K_B}\right\rbrace,
\end{split}
\end{align}
i.e.\ exactly the game formulation of the paper: $\xi_B$ is Alice's
assemblage, $N_B$ is Bob's complete multimeter, and the multimeter
normalization is the local equality constraint $F_B$ of the hierarchy. Four
experiments are run, distinguished by the local state space $K_B$ and one
optional constraint:

\begin{center}
\begin{tabular}{lll}
\toprule
experiment & $K_B$ & extra constraint\\
\midrule
classical bit & simplex $\Delta_2$ & --\\
square gbit (``boxworld'') & square, vertices $(1,\pm1,\pm1)$ & --\\
pentagon & regular pentagon, $r_5=\sqrt{\sec(\pi/5)}$ & --\\
pentagon + inequality & same & $q_{1\vert 0}\geq 0.7$, Eq.~\eqref{eq:GA}\\
\bottomrule
\end{tabular}
\end{center}

\noindent
Here $r_5$ denotes the pentagon's circumradius. With this value the
pentagon is \emph{strongly} self-dual: the extreme rays of the dual cone,
normalized by their first coordinate, coincide with the vertices of the
pentagon itself (odd polygons are strongly self-dual, even ones only weakly
so, the dual of the square being the square rotated by $\pi/4$ and
rescaled)~\cite{Janotta2011}; the notebook verifies both statements
numerically. The inequality experiment imposes the local constraint
\begin{align}\label{eq:GA}
    G_A(\xi_B)\;=\; q_{1\vert 0}(\xi_B) - 0.7\cdot 1_{\Kasm}(\xi_B)\;\geq\; 0,
\end{align}
forcing Alice to answer $a=1$ on question $x=0$ with probability
at least $0.7$; on the extension it is imposed in the conic sense
(item (2c) in the paper's definition of the hierarchy $\Sigma^n(A:B)$;
labels (2b)/(2c) below always refer to those constraint items, not to
equation numbers of this document), which is precisely the feature that
distinguishes the hierarchy of~\cite{paper} from the polarization
hierarchy.

{\sloppy
In the code, a state space is a \code{Cone} object (vertices \code{gens},
dual rays \code{dual}, order unit \code{u}, interior point \code{tau});
the derived game data live in a \code{GameProblem}
(objective matrix \code{W}, normalization \code{uA}, unit \code{t},
equality rows \code{F\_rows}, optional inequality rows \code{GA\_rows},
and the dual rays \code{DR\_A} and \code{DR\_B} of the two bodies). All
experiments are assembled in \code{make\_experiments()} of
\code{reproduce.py}.\par}

% ===========================================================================
\section{Every level is an exact linear program}\label{sec:lp}

For polytopic state spaces both bodies in~\eqref{eq:gamma} are polytopes.
Membership in the maximal tensor product of the corresponding cones admits
an exact finite facet description, because the dual of the maximal tensor
product is the minimal tensor product of the duals, whose extreme rays are
the products of the factors' extreme rays. Level $n$ of the hierarchy
relaxes the product $\xi_B\otimes N_B$ to a permutation-invariant,
normalized extension
\begin{align}\label{eq:leveln}
\begin{split}
    p^{(n)} = \max\quad & \big\langle W\otimes t^{\otimes(n-1)},\, Y\big\rangle\\
    \text{s.t.}\quad
    & Y \in \Casm\tmax \Cmm^{\tmax n}
      \qquad\qquad\qquad\quad\ \text{(facet inequalities, see below)}\\
    & \lrbracket{u_A\otimes t^{\otimes n}}(Y)=1
      \qquad\qquad\qquad\ \, \text{(normalization)}\\
    & \lrbracket{\operatorname{id}\otimes F_B\otimes \operatorname{id}^{\otimes(n-1)}}(Y)=0
      \qquad\ \text{(local equalities, (2b))}\\
    & \lrbracket{G_A\otimes \operatorname{id}^{\otimes n}}(Y)\geq 0
      \qquad\qquad\quad\ \, \text{(local inequalities, (2c))}
\end{split}
\end{align}
with the membership constraint written out as
\begin{align}\label{eq:facets}
    \lrbracket{\varphi\otimes\psi_{k_1}\otimes\cdots\otimes\psi_{k_n}}(Y)\;\geq\;0
    \qquad \text{for all }\varphi\in\ext \Casm^{\,*},\
    \psi_{k_i}\in\ext \Cmm^{\,*}.
\end{align}
Hence every level is a linear program (LP), and the LP optimum equals
$p^{(n)}$ exactly up to solver tolerances of order $10^{-8}$
(HiGHS~\cite{HiGHS}, dual simplex for the small levels, interior point with
crossover for the largest reduced LPs).

For the pentagon,
$\lrvert{\ext \Casm^{\,*}}=\lrvert{\ext \Cmm^{\,*}}=20$, so after the
multiset reduction of Section~\ref{sec:reductions} the family~\eqref{eq:facets}
contains $20\binom{19+n}{n}$ rows: $4\,200$ at $n=2$, about
$1.77\times 10^5$ at $n=4$, $8.5\times 10^5$ at $n=5$ and $3.5\times 10^6$
at $n=6$ (plus $\binom{19+n}{n}$ conic rows of $G_A$ in the constrained
experiment). The further reductions below compress these to LPs with a few
hundred variables that solve within seconds to minutes on one CPU core
(Section~\ref{sec:sizes}).

% ===========================================================================
\section{Three exact reductions}\label{sec:reductions}

All three reductions are \emph{exact}: they change neither the feasible
value set nor the optimum, only the parametrization. The paper does not
discuss them; Section~\ref{sec:unreduced} verifies numerically that the
levels solved with none of them give the same values.

\subsection{Permutation symmetry: the multiset basis}

By permutation invariance the extension $Y_{AB^n}$ is stored in the
symmetric (multiset) basis over the $n$ copies of $B$: one coordinate per
pair $(\alpha, M)$ of an $A$-index $\alpha$ and a multiset $M$ of size $n$
of $B$-indices, holding the common value of $Y_{\alpha, i_1\ldots i_n}$
for any ordering $(i_1,\ldots,i_n)$ of $M$ (``value form''). A symmetric
slot then has dimension $\binom{d_B+n-1}{n}$ instead of $d_B^{\,n}$, and a
one-slot contraction with a covector $f$ is the sparse update
\begin{align*}
    (f\mathbin{._1}Y)[\alpha, M'] \;=\; \sum_i f_i\, Y[\alpha, M'+\lbrace i\rbrace],
\end{align*}
implemented with a precomputed multiset-successor table (class
\code{SymSpace} of the library). The facet list~\eqref{eq:facets}
likewise reduces to ray \emph{multisets} $k_1\leq\cdots\leq k_n$.

\subsection{Local equalities as a subspace}

By permutation invariance, the equality constraint (2b) forces every
$B$-slot of $Y$ into $\ker F_B$, a subspace of dimension $1+2\dim K_B$.
Parametrizing all slots inside $\ker F_B$ from the start realizes
\emph{all} equality constraints exactly and removes them from the LP; the
single remaining equality row is the normalization. This also matters for
the rounding of Section~\ref{sec:rounding}: any point assembled from such
slots satisfies the multimeter normalization automatically.

\subsection{Dihedral symmetry: orbit averaging}\label{sec:orbit}

The full problem data of~\eqref{eq:leveln} --- objective, normalization,
both cones, $F_B$ and $G_A$ --- are invariant under the dihedral symmetry
group of the polygon $K_B$, acting simultaneously on all tensor factors.
Group elements permute the facets~\eqref{eq:facets} by relabelling
(a permutation of the dual rays inside $\varphi$ and of the vertex labels
inside the $\psi_{k}$; \code{dihedral\_actions()} constructs both
permutation actions). Replacing each facet by its orbit average therefore
yields an LP with the same optimum:
\begin{itemize}[leftmargin=1.5em]
\item averaged rows are convex combinations of valid rows, so the
averaged LP is a relaxation ($p^{(n)}_{\mathrm{avg}}\geq p^{(n)}$);
\item conversely, the group average of any feasible point of the averaged
LP is feasible for the original LP with the same objective value, since
objective and constraints are invariant ($p^{(n)}_{\mathrm{avg}}\leq p^{(n)}$).
\end{itemize}
The averaged rows span an invariant subspace of dimension
$45, 111, 246, 504$ for $n=2,\ldots,5$ (pentagon), and the LP is solved in
these coordinates. In symmetric coordinates rescaled by the square root of
the orbit multiplicity the group action is orthogonal, so the optimizer
reconstructed from the reduced solution is \emph{exactly} group-invariant
and hence satisfies every individual facet, which is what the rounding
step needs. Invariance does not mean uniqueness: the invariant optimal face
is itself degenerate in general, and the reconstructed optimizer is the one
the solver happens to return (Section~\ref{sec:sensitivity}).

\begin{note}\sloppy
The library contains two level solvers with the complete facet list:
\code{hierarchy\_level\_full} (multiset basis and $\ker F_B$ only, used
for the classical bit) and \code{hierarchy\_level\_symfull} (in addition
orbit-averaged; used for the square gbit and both pentagon experiments).
The \code{certified\_by} field in \code{results.json} records which one
produced each number. Both are complete-list LPs --- there is no
cutting-plane truncation anywhere in the reported numbers.
\end{note}

% ===========================================================================
\section{Feasible inner points: the rounding pipeline}\label{sec:rounding}

The inner sequence follows the constructive proof of the paper's
convergence theorem, applied to the optimizer $y^{(n)}$ of each level as
returned by the LP solver.

\begin{enumerate}[leftmargin=2em]
\item \textbf{Measurement.} The extreme rays $\psi_k$ of $\Cmm^{\,*}$
satisfy $\sum_k \psi_k = s\cdot 1_{\Kmm}$ for a scalar $s$ ($s=4, 8, 10$
for bit, square, pentagon), so $\lbrace \psi_k/s\rbrace_k$ is a symmetric
informationally complete measurement on $\Kmm$ (the code asserts this
property at runtime). For each $m\in\lbrace 0,\ldots,n-1\rbrace$, $m$ of the
$n$ multimeter copies of $y^{(n)}$ are measured with it.
\item \textbf{Conditioning.} For each outcome multiset $z^m$ (orderings
are equivalent by permutation invariance; the multinomial weight is
accounted for), the remaining $n-1-m$ copies are marginalized with the
unit $t$, leaving a conditional pair on $A\times B$.
\item \textbf{Productization.} The conditional pair is split into its
marginals $\lrbracket{x_{A\vert z^m}, x_{B\vert z^m}}$, i.e.\ a bona fide
product strategy (an assemblage and a multimeter).
\item \textbf{Report.} Iterating over all $m$ and all outcomes,
\begin{align}\label{eq:inner}
    L^{(n)}\;\coloneqq\;\max_{m}\ \max_{z^m}\
    \Val\lrbracket{x_{A\vert z^m}, x_{B\vert z^m}}\;\leq\;\gamma^*(G),
\end{align}
the maximization (sign-flipped) version of the theorem's construction,
appropriate for the game problems, which maximize a winning probability;
since each conditional product is feasible, the best of them is at least
as good as their mixture, the de Finetti inner point of the theorem (whose
value is recorded as well, field \code{avg\_at\_best\_m}).
\end{enumerate}

\noindent
\textbf{Why every candidate is feasible.} Cone membership survives
contraction with positive functionals (the $\psi_k$ and $t$ are positive);
the multimeter normalization $F_B$ holds because every slot is
parametrized inside $\ker F_B$ (Section~\ref{sec:reductions}); and the
local inequality $G_A$ survives conditioning because it is imposed on the
extension in the conic sense of constraint (2c). Hence each
$\lrbracket{x_{A\vert z^m}, x_{B\vert z^m}}$ is feasible for the
\emph{original, constrained} problem and $L^{(n)}$ is a lower bound
attained by an explicit strategy --- this is the feasibility mechanism of
the paper's proof, verified numerically: over all experiments and levels,
the worst feasibility slack of a rounded candidate in the shipped run is
$-8.4\times 10^{-11}$, i.e.\ all candidates are feasible to within
$10^{-10}$ (field \code{min\_slack} in \code{results.json}; the library
raises an error below $10^{-6}$). For the optimizers of the unreduced LPs
the same candidates are additionally certified \emph{rigorously}, in
interval arithmetic (Section~\ref{sec:sensitivity}).

Two practical remarks. First, the maximization over $m$
in~\eqref{eq:inner} matters: for every entry of Table~1 the best value is
attained at $m=n-1$ (for the classical bit among ties), and both the per-$m$
best and the per-$m$ mixture value are recorded. Second --- observation (e)
of the paper --- the optimal face of each level LP is typically degenerate,
the rounding map is nonlinear (conditioning divides by the outcome
probability, then productizes, then maximizes), and so the rounded value is
a property of the \emph{particular} optimizer that is rounded, not of the
level alone. Table~1 reports the values obtained from the optimizer
returned by our implementation in the run shipped as \code{results.json};
Section~\ref{sec:sensitivity} quantifies the spread.

% ===========================================================================
\section{Exact reference values}\label{sec:exact}

The optimal values themselves are computed directly in
\code{exact\_game\_value.py}, bypassing hierarchy and rounding. Both bodies
in~\eqref{eq:gamma} are polytopes and the objective is bilinear, so the
optimum is attained at a pair of vertices, and $\Kmm$ is affinely
isomorphic to $E(K_B)\times E(K_B)$, so its vertices are the pairs of
extreme effects.

\begin{enumerate}[leftmargin=2em]
\item \textbf{Extreme multimeters.} The vertices of the effect body
$E(K_B)=\lbrace e : 0\leq e(g_j)\leq 1\rbrace$ are enumerated exactly as
intersections of $d$ of the $2m$ planes $e(g_j)\in\lbrace 0,1\rbrace$: $4$
for the bit, $6$ for the square gbit, $12$ for the pentagon ($0$, $1_{K_B}$,
the five sharp effects $e_j$ with $e_j(g_j)=1$ and $e_j=0$ on the opposite
edge, and their complements $1_{K_B}-e_j$). Hence $16$, $36$ and $144$
multimeter pairs.
\item \textbf{One LP per multimeter pair.} For a fixed extreme multimeter
$N_B$ the remaining problem is the LP $\max\lbrace\langle W N_B,\xi\rangle :
\xi\in\Kasm\rbrace$ over the assemblage polytope ($3d$ variables,
$4m$ product facets, plus $G_A$ in the constrained experiment). The affine
symmetry group of $K_B$ (of order $2$, $8$, $10$) acts on the multimeter
pairs and leaves the LP value invariant, so only orbit representatives are
solved: $10$, $11$, $24$ of the $16$, $36$, $144$ pairs (\code{--full}
solves every pair and reproduces the same values).
\item \textbf{Exact certificates.} Each LP is solved in floating point
(HiGHS dual simplex without presolve, which returns a basic solution) and
then certified in exact arithmetic over $\mathbb{Q}(\sqrt5)$ ($\mathbb{Q}$
for bit and gbit): the optimal vertex is recomputed from its active facets
and verified feasible, and exact Farkas multipliers $y\geq 0$ with
$\mu\, u_A - \sum_j y_j\varphi_j = W N_B$ are recomputed from the LP basis,
so that $\mu=\langle W N_B,\xi^\star\rangle$ is the exact LP value. The
maximum over the certified values is the exact game value. The pentagon is
realized over $\mathbb{Q}(\sqrt5)$ as the affine image of the regular
pentagon under $\operatorname{diag}(1, 1/\sin 72^\circ)$ (game values are
affine invariants of $K_B$; the affine equivalence to the library's
pentagon is verified numerically to $10^{-15}$), and the float pipeline is
additionally re-run in the library's own coordinates
(\code{GameProblem.W}, facet lists), agreeing to $10^{-15}$.
\item \textbf{Vertex-enumeration cross-check.} Independently, all vertices
of $\Kasm$ ($8$, $24$, $60$, and $220$ with the constraint) are enumerated
by the double description method (\code{pycddlib}~\cite{cddlib}; exact GMP
arithmetic for the rational state spaces, floating point followed by exact
re-derivation of every vertex for the pentagon), and the maximum over
vertex pairs is recomputed exactly; it agrees with the LP value in every
case. (For the square gbit $\Kasm=CS^1_{2,2}\tmax K_B$ is an affine copy of
the $2$-$2$-$2$ no-signaling polytope, whence the $24$ vertices.)
\end{enumerate}

\noindent The results (CHSH, winning probability) are
\begin{align*}
 \gamma^*(\mathrm{CHSH}) &= \tfrac34,\qquad 1,\qquad \tfrac{2\sqrt5-1}{4}\approx 0.86803,\qquad
 \tfrac{3(9+\sqrt5)}{40}\approx 0.84271
\end{align*}
for the bit, the square gbit, the pentagon and the pentagon with the
constraint~\eqref{eq:GA}: the classical CHSH bound, the algebraic maximum
(PR-box value), and two closed forms in $\sqrt5$. The constraint binds
(observation (f)). For the pentagon the optimum is attained by $40$ of the
$144$ multimeter pairs; one optimal strategy in closed form is: Bob measures
the sharp effects $n_{1\vert 0}=e_3$, $n_{1\vert 1}=e_2$ (adjacent
vertices), and Alice's assemblage steers along the two diagonals through
their common crossing point, $\tilde\rho_{0\vert 0}=\varphi^{-1}g_0$,
$\tilde\rho_{1\vert 0}=\varphi^{-2}g_2$, $\tilde\rho_{0\vert 1}=\varphi^{-1}g_1$,
$\tilde\rho_{1\vert 1}=\varphi^{-2}g_4$ with $\varphi=(1+\sqrt5)/2$, so that
$\Bar\rho$ is the intersection of the diagonals $g_0g_2$ and $g_1g_4$ (which
it divides in the golden ratio) --- not the maximally mixed state. In the
constrained experiment the constraint is active, $q_{1\vert 0}=7/10$, and
$30$ pairs are optimal. The script runs in a few seconds and writes
\code{exact\_game\_values.json} (values, closed forms, maximizing
multimeters, exact dual certificates, vertex counts, comparison with
\code{results.json}) and \code{exact\_game\_values.md}.

% ===========================================================================
\section{Problem sizes and runtimes}\label{sec:sizes}

Table~\ref{tab:sizes} lists, for every experiment and level, the number of
LP variables (invariant coordinates after all reductions of
Section~\ref{sec:reductions}), the number of inequality rows of the
reduced LP, the wall-clock time of the outer solve, and the inner value
with its optimal $m$, all as recorded in the shipped \code{results.json}.
All computations ran on a single CPU core; total time for
\code{reproduce.py} is about $15$--$40$ minutes depending on the machine.

\begin{table}[ht]\centering\small
\begin{tabular}{llrrrrr}
\toprule
experiment & & $n=1$ & $n=2$ & $n=3$ & $n=4$ & $n=5$\\
\midrule
classical bit
 & variables            & 30 & 90 & 210 & 420 & 756\\
 & rows                 & 64 & 288 & 960 & 2640 & 6336\\
 & outer solve [s]      & 0.04 & 0.03 & 0.08 & 0.3 & 7.3\\
 & $L^{(n)}$ ($m^*$)    & 0.500 (0) & 0.750 (0) & 0.750 (2) & 0.750 (3) & 0.750 (2)\\
\midrule
square gbit
 & variables            & 15 & 45 & 123 & 282 & 606\\
 & rows                 & 32 & 288 & 1632 & 7824 & 31008\\
 & outer solve [s]      & 0.01 & 0.06 & 0.5 & 5.2 & 106\\
 & $L^{(n)}$ ($m^*$)    & 0.500 (0) & 0.563 (1) & 0.625 (2) & 1.000 (3) & 1.000 (4)\\
\midrule
pentagon
 & variables            & 15 & 45 & 111 & 246 & 504\\
 & rows                 & 48 & 456 & 3184 & 18012 & 85728\\
 & outer solve [s]      & 0.06 & 0.2 & 1.4 & 52 & 245\\
 & $L^{(n)}$ ($m^*$)    & 0.500 (0) & 0.546 (1) & 0.604 (2) & 0.707 (3) & 0.755 (4)\\
\midrule
pentagon + ineq.
 & variables            & 15 & 45 & 111 & 246 & 504\\
 & rows                 & 52 & 486 & 3364 & 18973 & 90160\\
 & outer solve [s]      & 0.01 & 0.1 & 1.3 & 20 & 293\\
 & $L^{(n)}$ ($m^*$)    & 0.511 (0) & 0.610 (1) & 0.612 (2) & 0.666 (3) & 0.740 (4)\\
\bottomrule
\end{tabular}
\caption{Reduced LP sizes, outer solve times and inner values,
as recorded in \code{results.json} (fields \code{n\_vars}, \code{n\_cuts},
\code{seconds}, \code{inner}). The classical bit is solved with the plain
full facet list (no orbit averaging); the other three experiments use the
dihedral orbit-averaged LP. For the pentagon the row counts realize the
$20\binom{19+n}{n}$ facets of Section~\ref{sec:lp} after multiset
reduction and orbit averaging. For the classical bit the $m^*$ in
parentheses is the first $m$ attaining the maximum; all $m$ tie from $n=2$
on. Times are from one run on a single core and vary between machines.}
\label{tab:sizes}
\end{table}

The symmetry reduction is generic for polygon models and keeps level
$n=6$ ($3.5\times 10^6$ facets before averaging) within reach of a laptop
(not benchmarked). For state spaces admitting conic lifts the same pipeline
transports to the lifting cones, with LPs replaced by the corresponding
conic programs.

% ===========================================================================
\section{Cross-check without any reduction}\label{sec:unreduced}

The paper states that every level is an exact LP and reports the values
without discussing the reductions of Section~\ref{sec:reductions}. To make
sure that the reductions are indeed exact in the implementation, the folder
\code{unreduced\_checks/} solves the \emph{same} level-$n$ programs with
none of them, exactly as written in the definition of $\Sigma^n(A:B)$:
\begin{itemize}[leftmargin=1.5em]
\item \textbf{Variables.} The full extension $Y\in\mathbb{R}^{d_A}\otimes
(\mathbb{R}^{d_B})^{\otimes n}$, one free real per tensor entry
($d_A=d_B=9$ for the polygon theories, $6$ for the classical bit); no
symmetric coordinates, no subspace.
\item \textbf{Constraints.} All product facets~\eqref{eq:facets} over
\emph{ordered} tuples ($20\cdot 20^n$ rows for the pentagon), permutation
invariance as explicit equality rows for the adjacent transpositions
(which generate $S_n$; invariance under all $n!$ permutations is checked a
posteriori), the normalization, the local equality (2b) on the first copy
(checked a posteriori on every copy), and in the fourth experiment the
conic rows (2c) over all ordered tuples.
\end{itemize}
The problem data ($W$, $u_A$, $t$, $F_B$, $G_A$, both extreme-ray
families) are rebuilt from the definitions in
\code{direct\_hierarchy.py} and checked to coincide with the library's
data, so both routes provably solve the same instance. Two formulations
are used:
\begin{enumerate}[leftmargin=2em,label=(\Alph*)]
\item \textbf{Explicit Kronecker matrix} (\code{direct\_hierarchy.py}):
every row is literally the Kronecker product of its defining vectors and
the whole matrix is handed to HiGHS. The facet matrix has
$\mathrm{nnz}(\Phi)\,\mathrm{nnz}(\Psi)^n$ nonzeros ($5.0\times10^7$ for
the pentagon at $n=3$, $4.2\times10^9$ at $n=4$), so this literal
formulation reaches $n\leq 3$ for the polygon theories and $n\leq 5$ for
the classical bit on a machine with a few GB of memory.
\item \textbf{Complete facet family with intermediate contractions.}
The same LP with the contraction spelled out slot by slot through
auxiliary variables
$Z_\varphi=(\varphi\otimes\mathrm{id}^{\otimes n})(Y)$,
$Z_{\varphi,k_1}=(\psi_{k_1}\otimes\mathrm{id}^{\otimes(n-1)})(Z_\varphi)$,
\ldots, and the facet rows $\psi_{k_n}(Z_{\varphi,k_1\ldots k_{n-1}})\geq 0$
--- the nested sums one would type into a modelling language
(\code{direct\_hierarchy\_tree.py}). Eliminating the $Z$'s returns the rows
of (A), so the LP has the same feasible $Y$'s and the same optimum; rows now have about five nonzeros, which brings
$n=4$ for the pentagon within reach ($5.8$ million rows, $2.6$ million
variables, $2.7\times 10^7$ nonzeros, about $2.6$ hours and $5$ GB).
\end{enumerate}
A result is recorded only if HiGHS reports model status \emph{Optimal}, and
the returned $Y$ is verified independently by tensor contraction (minimum
over all product facets and $G_A$ rows, normalization, $F_B$ on every copy,
invariance under all $n!$ permutations), with pass/fail flags at explicit
tolerances ($10^{-8}$). Every optimizer is stored with its SHA-256, the
solver options, iteration counts, residuals and the software environment
(\code{direct\_optimizers/}, \code{direct\_check\_runs/}).

\begin{table}[ht]\centering\small
\begin{tabular}{@{}llllr@{}}
\toprule
$K_B$ & $n$ & $p^{(n)}$ unreduced & $|$difference$|$ & formulation, time\\
\midrule
classical bit & 1--5 & 1, 0.75, 0.75, 0.75, 0.75 & $\leq 4.4\times10^{-16}$ & (A), $\leq 92$ s\\
square gbit & 1--3 & 1.5, 1, 1 & $\leq 1.7\times10^{-13}$ & (A), $\leq 42$ s\\
pentagon & 1--3 & 1.34164079, 0.92082039, 0.90420549 & $\leq 3.0\times10^{-14}$ & (A), $\leq 1012$ s\\
pentagon & 4 & 0.89488189 & $6.0\times10^{-14}$ & (B), 9485 s\\
pentagon + ineq. & 1--3 & 1.17331263, 0.91813655, 0.89521389 & $\leq 9.1\times10^{-13}$ & (A), $\leq 1529$ s\\
\bottomrule
\end{tabular}
\caption{Unreduced level values versus the reduced pipeline
(\code{unreduced\_checks/direct\_results\_table.md}; ten digits there). The
differences to \code{results.json} are at the level of the LP tolerances;
the reductions are exact. Formulation (B)
was additionally run at $n=2,3$ for the pentagon and at $n=2$ for the
constrained pentagon, agreeing with (A) to $10^{-14}$.}
\label{tab:unreduced}
\end{table}

Table~\ref{tab:unreduced} summarizes the outcome: for every entry of
Table~1 of the paper that is reachable without any reduction (all four
experiments at $n\leq 3$, the classical bit up to $n=5$, and via (B) the
pentagon at $n=4$), the unreduced value agrees with the reduced one to
within $10^{-12}$. This is the numerical confirmation that the three
reductions change the parametrization, not the optimum --- and it
quantifies why they are needed: at $n=4$ the literal facet matrix already
has $4\times 10^9$ nonzeros, at $n=5$ $3.5\times 10^{11}$.

% ===========================================================================
\section{Optimizer dependence of the inner values}\label{sec:sensitivity}

Observation (e) of the paper rests on the files of
\code{unreduced\_checks/} as well. The script
\code{direct\_inner\_rounding.py} carries out the rounding of
Section~\ref{sec:rounding} on the optimizers of the unreduced LPs, in the
full tensor picture (all ordered outcome tuples, all
assignments of the copies to the roles measured / kept / marginalized,
which agree to $10^{-14}$ by permutation invariance); every rounded
candidate is checked feasible in the unreduced description, the best
candidate is re-evaluated from its explicit states and effects, the same
optimizer mapped to the reduced basis is handed to the library's
\code{inner\_rounding} as an implementation cross-check (agreement to
$10^{-14}$), and the best candidate is turned into a \emph{rigorous} lower
bound by \code{direct\_inner\_certify.py}: the instance is rebuilt in
interval arithmetic (\code{mpmath}, 60 digits) from its definition, the
equality constraints are imposed exactly, the candidate is mixed with a
strictly interior product strategy ($\varepsilon=10^{-9}$), every
inequality of the original problem is verified strictly positive on the
mixed strategy, and the lower endpoint of the interval of its value is the
certified bound ($L^{(n)}-O(10^{-9})$; for the rational instances bit and
gbit the candidate is rationalized and checked exactly, giving $3/4$, $1$,
$11/18$, \ldots).

\code{direct\_inner\_sensitivity.py} then rounds, for $n=1,2$, the
optimizers of the same level LP returned by four solver routes --- the
explicit formulation (A) with dual simplex, with interior point plus
crossover, and with interior point without crossover, and the equivalent
formulation (B) with interior point --- each raw and after dihedral group
averaging. Table~\ref{tab:sens} shows the result at $n=2$: the outer value
is the same to $10^{-10}$ throughout while the optimizers differ by $O(0.1)$
in sup norm, and the rounded values spread over the ranges quoted in the
paper. Vertex optimizers (simplex, crossover) tended to round to higher
values than interior ones, but no route dominates; for the square gbit the
simplex vertex rounds to a strategy with PR-box correlations already at
$m=0$, so that for this optimizer the bracket closes at $n=2$ (the shipped
run closes it at $n=4$). Group averaging makes an optimizer invariant but
not unique: the averaged optimizers round to yet other values. At $n=1$
all routes give the same value ($1/2$, $1/2$, $1/2$ and
$(15-\sqrt5)/25$).

\begin{table}[ht]\centering\small
\begin{tabular}{lcccc@{\qquad}c}
\toprule
$K_B$ ($n=2$) & simplex & ipm\,+\,xover & ipm & tree, ipm & range\\
\midrule
classical bit & $0.7500$ & $0.7500$ & $0.5494$ & $0.5245$ & $[0.520,\ 0.750]$\\
square gbit & $1.0000$ & $0.6111$ & $0.5666$ & $0.5515$ & $[0.551,\ 1.000]$\\
pentagon & $0.5474$ & $0.5618$ & $0.5499$ & $0.5423$ & $[0.532,\ 0.562]$\\
pentagon + ineq. & $0.6136$ & $0.6146$ & $0.6080$ & $0.6065$ & $[0.606,\ 0.615]$\\
\bottomrule
\end{tabular}
\caption{Rounded values $L^{(2)}$ of the raw optimizers returned by four
solver routes for the same unreduced level-2 LP, and the range over the
routes including the group-averaged optimizers (the intervals quoted in
observation (e) of the paper). The shipped run (Table~1) has
$L^{(2)}=0.7500$, $0.5625$, $0.5459$, $0.6097$. Every value is certified
rigorously as a lower bound on $\gamma^*$; all data, hashes and certificates
are in \code{direct\_inner\_sensitivity.json}.}
\label{tab:sens}
\end{table}

The consequence for reproducibility is spelled out in the paper: $p^{(n)}$
is intrinsic to the level, whereas $L^{(n)}$ is attached to a specific
optimizer. HiGHS is deterministic only within one build and platform; the
same source with the same HiGHS version can return another vertex of the
optimal face on another machine. Concretely, re-running \code{reproduce.py} in a second environment
(Linux x86-64, HiGHS 1.15.1) reproduced all outer values to $5\times
10^{-11}$ and all inner values of Table~1 for the bit, the gbit and at
$n\leq 2$, but returned $L^{(3)}=0.6111$ for the pentagon instead of the
shipped $0.6044$, and values differing in the fourth decimal at $n=4,5$
for the pentagon and at $n=3,4$ for the constrained pentagon. This is why the optimizers behind the shipped inner
values of \code{unreduced\_checks/} are stored with their hashes and
rounded from file by default, and why the notebook checks the inner values
for validity rather than for digits. The flip side, also noted in the
paper, is a cheap practical lever: rounding costs little compared with the
LP, so one may round several optimizers and keep the best strategy.

% ===========================================================================
\section{Level-1 diagnostics}\label{sec:level1}

Observation (a) of the paper is produced by
\code{level1\_pseudocorrelations.py}: it re-solves the level-1 LPs (the
maximal tensor product $\Kasm\tmax\Kmm$) and contracts the optimizer with
the Choi tensor $\chi_{A(K_B)^\star}$ as in~\eqref{eq:gamma}, which yields a
normalized, no-signalling table $p(a,b\vert x,y)$. For the classical bit,
whose cone is simplicial, the table is the PR box and $p^{(1)}=1$ is the
no-signalling value of CHSH. For the square gbit the entries reach $-1/4$
and $p^{(1)}=3/2$; for the pentagon the minimum entry is $(1-3/\sqrt5)/2$
and $p^{(1)}=3/\sqrt5$; for the constrained pentagon $p^{(1)}=1.1733$ with
negative entries as well. Level $1$ is thus a valid but trivial bound, and
the non-trivial information enters with the symmetric-extension
constraints from $n=2$ on.

% ===========================================================================
\section{Reproducing the results}\label{sec:repro}

\paragraph{Dependencies.}
Python 3.10 or newer with \code{numpy}, \code{scipy}, \code{highspy}
(HiGHS~\cite{HiGHS}) and \code{matplotlib} (\code{requirements.txt});
additionally \code{nbformat}, \code{nbclient} and \code{ipykernel} to
rebuild the notebook, \code{pycddlib} $\geq 3.0$ (built against cddlib
with GMP) for the vertex-enumeration cross-check of
\code{exact\_game\_value.py} (\code{--no-vertices} runs without it), and
\code{mpmath} for the interval certificates of \code{unreduced\_checks/}
(\code{--no-certify} skips them). The shipped data were produced with
Python 3.11, NumPy 2.4, SciPy 1.17 and HiGHS 1.15.1.

\paragraph{Table 1 and Figure 1 from scratch.}
\begin{center}
\code{python3 reproduce.py}\qquad(\(\sim\)15--40 min, one core)
\end{center}
runs, per experiment, the outer levels $p^{(1)},\ldots,p^{(5)}$ and the
rounding at every level, checks the sandwich
$L^{(n)}\leq\gamma^*\leq p^{(n)}$ against the exact optima, writes
\code{results.json} (pass \code{--out} to keep the shipped file) and
renders \code{hierarchy\_convergence.pdf/.png} via \code{make\_figure.py};
\code{--nmax 3} gives a quick run in seconds.

\paragraph{Exact reference values.}
\begin{center}
\code{python3 exact\_game\_value.py}\qquad(seconds; \code{--full} solves every multimeter pair)
\end{center}
prints, per experiment, the exact $\gamma^*$ with closed form, the number
of certified LPs and of maximizing multimeter pairs, the vertex count of
$\Kasm$, the agreement with the library-coordinate float pipeline and with
\code{results.json}, and writes \code{exact\_game\_values.json/.md}.

\paragraph{Observation (a).}
\code{python3 level1\_pseudocorrelations.py} prints the level-1 tables
(Section~\ref{sec:level1}).

\paragraph{Unreduced cross-check and observation (e)} (folder
\code{unreduced\_checks/}; all paths are resolved relative to the scripts,
so the commands work from any directory):
\begin{center}\small
\begin{tabular}{@{}l@{\quad}>{\raggedright\arraybackslash}p{58mm}@{}}
\code{python3 direct\_hierarchy.py --sizes} & explicit LP sizes per level\\
\code{python3 direct\_hierarchy.py} & (A), all levels that fit ($\sim$1 h); saves the optimizers\\
\code{python3 direct\_hierarchy\_tree.py --exp pentagon} & \\
\code{\ \ --n 4 --solver ipm --crossover off --ipm-tol 1e-10} & (B), pentagon $n=4$ ($\sim$2.6 h, $\sim$5 GB)\\
\code{python3 direct\_summary.py} & merges the run files into \code{direct\_results.json}\\
\code{python3 direct\_inner\_rounding.py --nmax 3} & rounds the saved optimizers, certifies (seconds)\\
\code{python3 direct\_inner\_sensitivity.py} & Table~\ref{tab:sens} from the saved route optimizers\\
\code{python3 direct\_inner\_sensitivity.py --resolve} & the four routes afresh on this machine\\
\end{tabular}
\end{center}

\noindent The last four commands regenerate the JSON files and Markdown
tables shipped in the folder from the stored run files and optimizers;
re-solving (\code{--resolve}, or running \code{direct\_hierarchy.py})
overwrites the stored optimizers and may change the rounded values, as
explained in Section~\ref{sec:sensitivity}.

\paragraph{Notebook.}
\code{reproduce\_results.ipynb} is a pre-executed walkthrough of all of the
above with $95$ soft assertion checks (levels $n\leq 3$ live, $n=4,5$ from
\code{results.json}; \code{RUN\_FULL = True} recomputes them);
\code{python3 build\_notebook.py --execute} regenerates and re-executes it
in about a minute.

\paragraph{What to test after modifying the pipeline.}
Test \emph{validity} and \emph{trends}, not digits of the inner values:
outer values must reproduce to $\sim 10^{-12}$ and be nonincreasing in
$n$ and $\geq\gamma^*$; rounded values must be $\leq\gamma^*$ with
feasibility slacks $\geq -10^{-10}$; their digits are optimizer dependent
(Section~\ref{sec:sensitivity}).

% ===========================================================================
\section{File-by-file guide}\label{sec:files}

\begin{center}\small
\begin{tabular}{>{\raggedright\arraybackslash}p{52mm}>{\raggedright\arraybackslash}p{90mm}}
\toprule
file & contents\\
\midrule
\code{README.md} & Overview, installation, the table, how to reproduce.\\
\code{implementation\_notes.tex/.pdf} & This document.\\
\code{gpt\_definetti\_games.py} & Library: polygon theories and their
cones (\code{Cone}, \code{pentagon\_selfdual}, \code{square\_gbit},
\code{classical\_bit}), game bodies (\code{GameProblem}), the two
complete-list level solvers (\code{hierarchy\_level\_full},
\code{hierarchy\_level\_symfull}), symmetric-tensor machinery
(\code{SymSpace}), \code{inner\_rounding}.\\
\code{reproduce.py} & Headless reproduction of Table 1 and the figure;
writes \code{results.json}.\\
\code{results.json} & The shipped run: per experiment the exact optimum,
per level the outer value with LP statistics (\code{n\_vars},
\code{n\_cuts}, \code{seconds}, \code{certified\_by}) and the rounding
(\code{best}, \code{best\_m}, per-$m$ values, \code{min\_slack}).\\
\code{make\_figure.py} & Renders \code{hierarchy\_convergence.pdf/.png}
from \code{results.json}.\\
\code{exact\_game\_value.py} & Exact direct computation of $\gamma^*$
(Section~\ref{sec:exact}); output \code{exact\_game\_values.json/.md}.\\
\code{level1\_pseudocorrelations.py} & Level-1 tables of observation (a).\\
\code{reproduce\_results.ipynb} & Pre-executed walkthrough ($95$ checks);
regenerated by \code{build\_notebook.py}.\\
\code{requirements.txt} & Python dependencies.\\
\midrule
\multicolumn{2}{l}{\emph{folder} \code{unreduced\_checks/}}\\
\code{direct\_hierarchy.py} & Formulation (A): the
literal level-$n$ LP; saves optimizers to \code{direct\_optimizers/}.\\
\code{direct\_hierarchy\_tree.py} & Formulation (B).\\
\code{direct\_summary.py} & Merges the run records of
\code{direct\_check\_runs/} into \code{direct\_results.json} and
\code{direct\_results\_table.md}.\\
\code{direct\_inner\_rounding.py} & Rounding of the
unreduced optimizers with all cross-checks; output
\code{direct\_inner\_results.json}, \code{direct\_inner\_table.md}.\\
\code{direct\_inner\_certify.py} & Rigorous interval /
exact-rational certificates of rounded strategies.\\
\code{direct\_inner\_sensitivity.py} & Solver-route
study of Section~\ref{sec:sensitivity}; output
\code{direct\_inner\_sensitivity.json/.md}.\\
\code{direct\_optimizers/} & Stored optimizers
(\code{*\_Y.npy}) with side-car JSON (SHA-256, solver options, residuals,
environment).\\
\code{direct\_check\_runs/} & Raw run records of (A) and
(B), including the pentagon $n=4$ tensor.\\
\bottomrule
\end{tabular}
\end{center}

\section*{Acknowledgements}
We thank Tobias Rippchen for fruitful discussions. This document and the
underlying code were prepared with the assistance of Anthropic's Claude
Fable 5 and OpenAI's ChatGPT 5.6.

% ===========================================================================
\begin{thebibliography}{9}

\bibitem{paper}
J.~A. Zeiss, G.~Kossmann, R.~Schwonnek, M.~Pl\'avala,
\emph{Finite de Finetti for convex bodies and polynomial optimization},
arXiv:2601.15184.

\bibitem{HiGHS}
Q.~Huangfu, J.~A.~J. Hall,
\emph{Parallelizing the dual revised simplex method},
Mathematical Programming Computation \textbf{10}, 119--142 (2018).
doi:\href{https://doi.org/10.1007/s12532-017-0130-5}{10.1007/s12532-017-0130-5}.

\bibitem{Janotta2011}
P.~Janotta, C.~Gogolin, J.~Barrett, N.~Brunner,
\emph{Limits on nonlocal correlations from the structure of the local state space},
New Journal of Physics \textbf{13}, 063024 (2011).
doi:\href{https://doi.org/10.1088/1367-2630/13/6/063024}{10.1088/1367-2630/13/6/063024}.

\bibitem{cddlib}
K.~Fukuda, A.~Prodon,
\emph{Double description method revisited},
in: Combinatorics and Computer Science, Lecture Notes in Computer Science
\textbf{1120}, Springer, 91--111 (1996).

\end{thebibliography}

\end{document}
